% ===== PRÉAMBULE CANONIQUE — à coller VERBATIM en tête de CHAQUE .tex =====
\documentclass[11pt,a4paper]{article}
\usepackage[utf8]{inputenc}\usepackage[T1]{fontenc}\usepackage{lmodern}
\usepackage[french]{babel}
\usepackage{amsmath,amssymb,mathtools}\usepackage{icomma}
\usepackage[version=4]{mhchem}          % \ce{...} : équations & formules
\usepackage{siunitx}\sisetup{locale=FR} % \qty{}{}, \si{}, \ang{}
\usepackage{chemfig}                    % \chemfig{...} : structures développées
\usepackage{xcolor}\usepackage{colortbl}
\usepackage{tikz}\usepackage{pgfplots}\pgfplotsset{compat=1.18}
\usepackage[most]{tcolorbox}\usepackage{enumitem}\usepackage{array,booktabs}
\usepackage[a4paper,margin=2.2cm]{geometry}
\usetikzlibrary{arrows.meta,calc,angles,quotes,positioning,patterns,backgrounds,shapes.geometric}
\definecolor{primary}{RGB}{222,49,99}\definecolor{primarydark}{RGB}{150,22,55}
\definecolor{defcol}{RGB}{0,114,178}\definecolor{methcol}{RGB}{52,92,148}
\definecolor{excol}{RGB}{0,135,90}\definecolor{attncol}{RGB}{200,40,55}
\tcbset{boxbase/.style={enhanced,breakable,boxrule=0.9pt,arc=2.5pt,
  left=3mm,right=3mm,top=2.3mm,bottom=2.3mm,fonttitle=\bfseries,coltitle=white}}
% --- boîtes TUTORIEL (noms EXACTS, ne pas renommer) ---
\newtcolorbox{cle}[1][]{boxbase,colback=defcol!6,colframe=defcol,
  title={\textsc{À retenir}\ifx\relax#1\relax\relax\else\textmd{ : \itshape #1}\fi}}
\newtcolorbox{methode}[1][]{boxbase,colback=methcol!7,colframe=methcol,sharp corners,
  title={\textsc{Méthode}\ifx\relax#1\relax\relax\else\textmd{ --- \itshape #1}\fi}}
\newtcolorbox{exemplebox}[1][]{boxbase,colback=excol!5,colframe=excol,
  title={\textsc{Exemple}\ifx\relax#1\relax\relax\else\textmd{ : \itshape #1}\fi}}
\newtcolorbox{piege}[1][]{boxbase,colback=attncol!6,colframe=attncol,
  title={\textsc{Piège}\ifx\relax#1\relax\relax\else\textmd{ : \itshape #1}\fi}}
\newtcolorbox{remarque}[1][]{boxbase,colback=black!4,colframe=black!45,
  title={\textsc{Remarque}\ifx\relax#1\relax\relax\else\textmd{ : \itshape #1}\fi}}
% --- boîtes EXERCICE / CORRIGÉ (noms EXACTS, auto-comptées) ---
\newtcolorbox[auto counter]{exo}[1][]{boxbase,colback=primary!4,colframe=primary,
  title={\textsc{Exercice}~\thetcbcounter\ifx\relax#1\relax\relax\else\textmd{ --- \itshape #1}\fi}}
\newtcolorbox[auto counter]{corrige}[1][]{boxbase,colback=excol!4,colframe=excol,
  title={\textsc{Corrigé}~\thetcbcounter\ifx\relax#1\relax\relax\else\textmd{ --- \itshape #1}\fi}}
\newcommand{\R}{\mathbb{R}}\newcommand{\N}{\mathbb{N}}\newcommand{\Z}{\mathbb{Z}}\newcommand{\ds}{\displaystyle}
% (un \usepackage{fancyhdr}+en-tête est possible ; il est ignoré côté web)
% =========================================================================
\begin{document}

\begin{exo}[l'échelle des oxacides du chlore]
Ces quatre oxacides du chlore diffèrent par le n.o.\ du chlore. Nomme chacun (préfixe/suffixe traditionnel).
\begin{enumerate}[label=\alph*)]
  \item \ce{HClO} (chlore $+\mathrm{I}$)
  \item \ce{HClO2} (chlore $+\mathrm{III}$)
  \item \ce{HClO3} (chlore $+\mathrm{V}$)
  \item \ce{HClO4} (chlore $+\mathrm{VII}$)
\end{enumerate}
\end{exo}

\begin{exo}[de l'oxacide à son anion : nom et formule]
Donne le \emph{nom} et la \emph{formule} de l'anion obtenu par perte des \ce{H+}.
\begin{enumerate}[label=\alph*)]
  \item acide hypochloreux \ce{HClO}
  \item acide chlorique \ce{HClO3}
  \item acide perchlorique \ce{HClO4}
\end{enumerate}
\end{exo}

\begin{exo}[de l'anion à l'oxacide]
Remonte de l'anion à l'acide : donne le \emph{nom} de l'acide et sa \emph{formule}.
\begin{enumerate}[label=\alph*)]
  \item nitrite \ce{NO2-}
  \item sulfate \ce{SO4^2-}
  \item carbonate \ce{CO3^2-}
\end{enumerate}
\end{exo}

\begin{exo}[transposer l'échelle au brome]
La même échelle hypo-/-eux/-ique/per- s'applique au brome. Nomme ces oxacides.
\begin{enumerate}[label=\alph*)]
  \item \ce{HBrO}
  \item \ce{HBrO3}
  \item \ce{HBrO4}
\end{enumerate}
\end{exo}

\begin{exo}[distinguer -ite et -ate : écrire la formule de l'anion]
Écris la formule (avec la charge) de chaque anion.
\begin{enumerate}[label=\alph*)]
  \item sulfite
  \item sulfate
  \item nitrite
  \item nitrate
\end{enumerate}
\end{exo}

\end{document}
